\section{Instance-Level Semantics for Constraint-Aware Counterfactual Queries}

% \input{section/4_instance/table/modality}
\label{sec:instance}

Given a query $q$, counterfactual retrieval is governed by instance-level semantics: candidates must satisfy feasibility constraints while remaining competitive under the transformation cost $c_\omega$. However, how these requirements are enforced and where they become computationally challenging varies across data modalities. In this section, we examine how P1 (traceability), P2 (cost faithfulness), and P3 (efficiency) manifest across modalities, focusing on where admissibility can be certified, how costs align with meaningful edits, and how constraint evaluation interacts with system design. The remaining set-level principle, P4, is discussed separately in Section~\ref{sec:set}.

\paragraph{Tabular and Relational Data.}
Tabular data offers the clearest setting for counterfactual querying: instance-level semantics are explicit, and edits correspond directly to feature changes. A counterfactual $\Delta(x_q,x)$ is simply a set of column-wise modifications, and the cost $\omega(x_q,x)$ can naturally aggregate per-attribute effort. Feasibility is also transparent—immutable attributes can be enforced directly by assigning infinite cost. As a result, tabular data aligns cleanly with traceability and cost faithfulness (P1 and P2), making it a natural starting point for data-centric counterfactual analysis.

However, this apparent simplicity is deceptive. The key challenge is that counterfactual retrieval is not purely geometric, even in tabular spaces. First, edit costs are heterogeneous: attributes differ in scale, semantics, and actionability, so $\omega$ must encode domain meaning rather than rely on naive distance metrics. Second, constraints induce structured partitions of the data space, where immutable or categorical rules define admissible regions. Third, ranking and feasibility are fundamentally decoupled: proximity in feature space does not guarantee admissibility, and ignoring constraints leads to wasted computation on infeasible candidates. These observations highlight that counterfactual querying is inherently a constraint-driven retrieval problem rather than a nearest-neighbor task. Efficient systems should therefore integrate column-level semantics into indexing and search, enforcing feasibility early and focusing computation on viable alternatives (P3).

\paragraph{Sequence and Temporal Data.}
In sequential domains, counterfactuals operate over trajectories rather than independent features. A query $x_q=(x_{q,1},\dots,x_{q,T})$ must be transformed into another valid sequence, so measuring cost $\omega(x_q,x)$ typically relies on alignment-based similarity, such as DTW~\cite{senin2008dynamic} or learned embeddings~\cite{nalmpantis2019signal2vec}. However, unlike tabular data, similarity alone is insufficient: temporal data is governed by dynamics, and admissibility depends on whether a sequence can be executed as a valid trajectory. This leads to a fundamental separation: cost is computed in observation space, while feasibility is determined in execution space. Two sequences may appear similar under alignment, yet correspond to invalid or unreachable states (e.g., illegal move sequences in chess). As a result, counterfactual retrieval becomes a problem of validating trajectories, not just comparing them.

This gap introduces three key challenges. First, edits must be \emph{execution-aware}: $\Delta(x_q,x)$ cannot be expressed as independent changes, but must correspond to coherent transformations under system dynamics, often requiring trajectory-level alignment or transport-based mappings (P1). Second, cost is \emph{temporally sensitive}: small timing shifts can affect $\omega(x_q,x)$, making it difficult to align similarity with intervention effort (P2). Third, sequences exhibit \emph{path dependency}: early changes propagate and alter downstream states, so admissibility must be verified through execution rather than local checks, shifting computation toward dynamic validation (P3). These challenges highlight that counterfactual querying in temporal data is inherently dynamics-aware.

\paragraph{Graph and Network Data.}

In graph domains, instances are structured objects with nodes and edges, and counterfactuals correspond to structural edits $\Delta(x_q,x)$. While embeddings $\phi(\cdot)$ enable scalable similarity search (e.g., $\|\phi(x_q)-\phi(x)\|_2$), they do not capture structural validity. Admissibility is instead governed by constraints in $\omega$, which enforce invariants such as connectivity, typing, or capacity. As in temporal data, cost and feasibility reside in different semantic layers: similarity is geometric, but validity is structural.

This mismatch introduces three challenges. First, graphs exhibit \emph{multi-level semantics}: edits can occur at the level of nodes, edges, or subgraphs, and small local changes may induce large structural effects, complicating traceable edits (P1). Second, there is \emph{embedding–constraint misalignment}: candidates close in embedding space may violate structural invariants, weakening the link between cost and feasible transformations (P2). Third, \emph{invariant verification is expensive}: many constraints are global, requiring whole-graph checks that can dominate query time (P3). These challenges suggest that graph counterfactuals cannot rely on embedding search alone. Effective systems must integrate structural constraints into retrieval, for example by restricting admissible edit spaces, maintaining summaries of invariants, or combining embedding search with constraint-aware validation.

\paragraph{Text.}
In text, counterfactuals correspond to edits $\Delta(x_q,x)$ over token sequences that carry rich structure, including syntax, discourse, and factual semantics. While pretrained embeddings~\cite{devlin2019bert} enable scalable similarity search, surface or embedding proximity alone cannot capture semantic equivalence. Admissibility is instead defined by constraints in $\omega$, which require edits to preserve coherence, entailment, and factual correctness. Moreover, feasible changes may occur at different semantic levels, such as topics, tone, sentiment, or keywords, making the formulation of counterfactual edits nontrivial. Therefore, cost and feasibility lie in different semantic layers: similarity is distributional, while validity is semantic.

This gap introduces three challenges. First, \emph{granularity mismatch}: edits are local (tokens or phrases), but their effects are global, potentially altering intent, discourse, or factual meaning, making $\Delta(x_q,x)$ hard to interpret (P1). Second, \emph{similarity--validity misalignment}: texts close in embedding space may violate entailment or factual constraints, weakening the connection between cost and meaningful change (P2). Third, \emph{context-dependent verification}: feasibility often cannot be checked locally and may require full-sequence reasoning or external knowledge, turning validation into a global and expensive process (P3).

\paragraph{Image.}
Image data presents the greatest challenge for counterfactual querying. An image is a high-dimensional pixel array mapped into an embedding space for scalable retrieval, yet pixel-level changes do not correspond directly to semantic edits. Small perturbations may cause large embedding shifts without meaningful object-level changes, while subtle visual differences can alter task-critical attributes~\cite{narodytska2017simple}. Feasibility, however, is defined at the level of objects, regions, or scene structure. Thus, admissibility under $\omega$ must enforce object-level transformations that satisfy domain constraints such as anatomical validity or physical plausibility~\cite{nam2026causal}, separating similarity from semantic validity.

This mismatch introduces three challenges. First, \emph{semantic opacity}: pixel differences do not translate into interpretable edits, making it difficult to construct a meaningful $\Delta(x_q,x)$ (P1). Second, \emph{perceptual–task inconsistency}: embedding similarity reflects visual closeness but may ignore task-relevant regions, weakening the alignment between cost and actionable change (P2). Third, \emph{region-aware verification}: feasibility often depends on localization or scene understanding, requiring additional vision modules that can dominate runtime (P3). These challenges highlight that image counterfactuals require an intermediate semantic layer. Approaches such as concept bottleneck models~\cite{de2025causally,koh2020concept,dominici2025counterfactual} map images to interpretable concepts, enabling cost and constraints to operate on meaningful attributes.


\paragraph{Extended to Multi-Modality.}
When data spans multiple modalities, counterfactual querying becomes a coordinated validation problem. Feasibility is defined jointly by modality-specific constraints and cross-modal consistency requirements~\cite{chen2019cross}. A valid candidate must satisfy the structural or semantic rules of each modality while remaining coherent across them (e.g., text, images, and attributes must agree). As modalities accumulate, both admissibility and cost modeling become more complex. In particular, a single semantic dimension (e.g., sentiment, object identity, or intent) may be expressed across multiple modalities, making it challenging to quantify transformation cost $\omega(x_q,x)$ in a fair and intrinsic way without double counting or inconsistency. At the same time, admissibility complexity can grow combinatorially.

In this setting, similarity search—whether in joint or modality-specific embeddings—serves mainly to propose candidates. Correctness is determined by cross-modal validation, which enforces both intra-modal constraints and inter-modal consistency. This shift highlights a key insight: \emph{multi-modal counterfactual querying is fundamentally verification-driven}. The central challenge is no longer retrieving similar instances, but designing principled methods to jointly evaluate feasibility and cost across modalities.